\chapter{Matrix Lie Groups}

\begin{center}\itshape
Rotations, unitary matrices and matrices of determinant one form groups that are also smooth manifolds, and the two structures lock together. Everything
about such a group near its identity is captured by a vector space with a bracket, its Lie algebra, and the exponential carries that linear world onto the
curved group. The prize at the end of this chapter is a fact about the world: a rotation by $360^\circ$ is not the same as no rotation at all, although it
looks the same, and the electron knows it.
\end{center}

In the summer of 1925, two graduate students of Paul Ehrenfest in Leiden, George Uhlenbeck and Samuel Goudsmit, were trying to make sense of the
spectra of alkali atoms, whose lines split in pairs. Their proposal was that the electron itself spins, carrying an angular momentum of half the
quantum unit and a magnetic moment. Ehrenfest liked it and suggested they ask Hendrik Lorentz, the grand old man of electron theory. Lorentz came back
some weeks later with pages of calculations: an electron of the classical size, spinning that fast, would have its surface moving at about ten times
the speed of light. The two students asked Ehrenfest to withdraw the paper. He answered that he had already sent it to \emph{Die Naturwissenschaften},
and, as Uhlenbeck remembered it fifty years later, that they were both young enough to be able to afford a stupidity.

It was not a stupidity, but it was strange. The spin of the electron is not a rotation of a small ball: when Wolfgang Pauli wrote down its quantum
mechanics in 1927, the states of a spinning electron formed a two-dimensional complex space on which rotations act through the group $SU(2)$ of
unitary $2\times2$ matrices, not through the rotation group $SO(3)$. The two groups look the same near the identity, but $SU(2)$ is a double cover:
each rotation corresponds to two matrices $\pm U$. A rotation by $2\pi$ brings an object back to itself, but it multiplies an electron's state by $-1$;
only a rotation by $4\pi$ is the identity. In 1975 two groups of physicists, one in Vienna and one in Missouri, split a beam of neutrons, turned the spin of
one half by a magnetic field, recombined the beams, and saw the interference pattern flip sign at $2\pi$ and return at $4\pi$.

The mathematics of this is the subject of this chapter. Matrix groups are manifolds (the theorem of von Neumann and Cartan); their tangent
space at the identity is a Lie algebra; continuous homomorphisms become linear maps; and the map $SU(2)\to SO(3)$ is the universal covering of the rotation
group, whose fundamental group is $\Z/2$. That is the plate trick of Tome~2, now proved.

\section{Matrix groups and their Lie algebras}

We regard $GL_n(\C)$ as a subset of $GL_{2n}(\R)$ (a complex number $a+ib$ acts on $\R^2$ by the matrix with rows $(a,-b)$ and $(b,a)$), and
$M_n(\C)$ as a real vector space. Fix the norm $\|X\|=\sup_{|v|=1}|Xv|$, which satisfies $\|XY\|\le\|X\|\|Y\|$.

\begin{definition}
A \emph{matrix Lie group} is a subgroup $G$ of $GL_n(\R)$ or $GL_n(\C)$ that is closed in $GL_n$. Its \emph{Lie algebra} is
\[\mathfrak g=\{X\in M_n:e^{tX}\in G\text{ for all }t\in\R\}.\]
\end{definition}

Examples: $SL_n(\R)$ ($\det=1$), $O(n)$ ($A^{\mathsf T}A=I$), $SO(n)=O(n)\cap SL_n(\R)$, $U(n)$ ($A^*A=I$), $SU(n)=U(n)\cap SL_n(\C)$, the Heisenberg group of upper triangular $3\times3$ real matrices with ones
on the diagonal, and every finite group of matrices. Each is defined by continuous equations, hence closed. For $\alpha$ irrational, the line $\{(e^{2\pi it},e^{2\pi i\alpha t}):t\in\R\}$ is a subgroup of the torus $U(1)\times U(1)$
that is dense and not closed (Figure~\ref{fig:torusflow} shows such a line), and the theory below does not apply to it.

The exponential $e^X=\sum X^k/k!$ was studied in Tome~S1, Proposition~3.13: $e^{X+Y}=e^Xe^Y$ when $XY=YX$, $\det e^X=e^{\operatorname{tr}X}$, $\frac{\dd}{\dd t}e^{tX}=Xe^{tX}$, and $e^{PXP^{-1}}=Pe^XP^{-1}$.
Near the identity it has an inverse, the logarithm
\[\log A=\sum_{k\ge1}\frac{(-1)^{k+1}}k(A-I)^k\qquad(\|A-I\|<1).\]

\begin{lemma}\label{lem:log}
\begin{enumerate}
\item[(a)] For $\|A-I\|<1$, $e^{\log A}=A$; for $\|X\|<\log2$, $\|e^X-I\|<1$ and $\log e^X=X$.
\item[(b)] There is a constant $c$ with $\|\log A-(A-I)\|\le c\|A-I\|^2$ for $\|A-I\|\le\frac12$.
\item[(c)] $\exp$ maps a neighbourhood of $0$ in $M_n$ diffeomorphically onto a neighbourhood of $I$, with differential the identity at $0$.
\end{enumerate}
\end{lemma}

\begin{proof}
(a) For a complex number, $e^{\log(1+z)}=1+z$ when $|z|<1$ and $\log e^w=w$ when $|w|<\log2$; these are identities between absolutely convergent power series, so they hold as identities of formal
power series after rearrangement, and therefore when the scalar is replaced by a matrix of norm in the same range (the rearrangements are justified by absolute convergence of
$\sum\|X\|^k/k!$ and $\sum\|A-I\|^k$). The bound $\|e^X-I\|\le e^{\|X\|}-1<1$ gives the first claim of the second part. (b) $\sum_{k\ge2}\frac1k\|A-I\|^k\le\|A-I\|^2\sum_{k\ge0}2^{-k}$. (c) $\exp$ is smooth (a power
series), and $e^{X}=I+X+O(\|X\|^2)$; the inverse function theorem (Theorem~\ref{thm:inverse}) applies, and (a) identifies the inverse with $\log$.
\end{proof}

\begin{proposition}\label{prop:oneparameter}
Every continuous homomorphism $\gamma\colon\R\to GL_n$ is $\gamma(t)=e^{tX}$ for a unique $X$, namely $X=\gamma'(0)$.
\end{proposition}

\begin{proof}
Choose $\varepsilon>0$ with $\|\gamma(t)-I\|<\frac12$ and $\|\log\gamma(t)\|<\frac12\log2$ for $|t|\le\varepsilon$, and let $X=\frac1\varepsilon\log\gamma(\varepsilon)$. The matrix $B=\log\gamma(\frac\varepsilon2)$ has $e^{2B}=\gamma(\varepsilon)=e^{\varepsilon X}$
with $\|2B\|<\log2$ and $\|\varepsilon X\|<\log2$, so by (a), $2B=\varepsilon X$, that is $\gamma(\frac\varepsilon2)=e^{\frac\varepsilon2X}$. By induction $\gamma(\varepsilon/2^k)=e^{\varepsilon X/2^k}$, then by the homomorphism property
$\gamma(m\varepsilon/2^k)=e^{m\varepsilon X/2^k}$ for all integers $m$, and by continuity and density $\gamma(t)=e^{tX}$ for all $t$. Differentiating at $0$ gives $X=\gamma'(0)$, hence uniqueness.
\end{proof}

\begin{proposition}\label{prop:liealgebra}
For $X,Y\in M_n$, with $[X,Y]=XY-YX$,
\[e^{X+Y}=\lim_{m\to\infty}\bigl(e^{X/m}e^{Y/m}\bigr)^m\quad\text{(Lie's product formula)},\]
\[e^{[X,Y]}=\lim_{m\to\infty}\bigl(e^{X/m}e^{Y/m}e^{-X/m}e^{-Y/m}\bigr)^{m^2}.\]
Consequently the Lie algebra $\mathfrak g$ of a matrix Lie group is a real vector space closed under the bracket, and $gXg^{-1}\in\mathfrak g$ for $g\in G$, $X\in\mathfrak g$.
\end{proposition}

\begin{proof}
Expanding, $e^{X/m}e^{Y/m}=I+\frac{X+Y}m+O(\frac1{m^2})$, so by Lemma~\ref{lem:log}(b) its logarithm is $Z_m=\frac{X+Y}m+O(\frac1{m^2})$, and $(e^{X/m}e^{Y/m})^m=e^{mZ_m}$ with
$mZ_m=X+Y+O(\frac1m)$. For the second formula, expand to second order: with $s=\frac1m$,
\[e^{sX}e^{sY}e^{-sX}e^{-sY}=I+s^2(XY-YX)+O(s^3),\]
because the terms of order $s$ cancel and those of order $s^2$ are $\frac{X^2}2+\frac{Y^2}2+\frac{X^2}2+\frac{Y^2}2+XY-X^2-XY-YX-Y^2+XY=XY-YX$. Its logarithm is $s^2[X,Y]+O(s^3)$, and the $m^2$-th power is
$e^{[X,Y]+O(1/m)}$.

If $X,Y\in\mathfrak g$, the products on the right lie in $G$ for every $t$ (replace $X,Y$ by $tX,tY$), and $G$ is closed, so $e^{t(X+Y)}$ and $e^{t^2[X,Y]}$ lie in $G$; for negative multiples of $[X,Y]$ use
$e^{-[X,Y]}=e^{[Y,X]}$. Clearly $sX\in\mathfrak g$ for real $s$. Finally $e^{tgXg^{-1}}=ge^{tX}g^{-1}\in G$.
\end{proof}

A real vector space with a bilinear, antisymmetric bracket satisfying the Jacobi identity $[X,[Y,Z]]+[Y,[Z,X]]+[Z,[X,Y]]=0$ is a \emph{Lie algebra}; the commutator of matrices satisfies it (expand). The
Lie algebras of the examples follow from the definition, by differentiating the defining equation along $e^{tX}$ at $t=0$ and conversely integrating:
\[\mathfrak{sl}_n=\{\operatorname{tr}X=0\},\quad\mathfrak{so}(n)=\{X^{\mathsf T}=-X\},\quad\mathfrak u(n)=\{X^*=-X\},\quad\mathfrak{su}(n)=\mathfrak u(n)\cap\mathfrak{sl}_n(\C).\]
For instance $\det e^{tX}=e^{t\operatorname{tr}X}$ is $1$ for all $t$ exactly when $\operatorname{tr}X=0$, and $(e^{tX})^*e^{tX}=e^{tX^*}e^{tX}$ is $I$ for all $t$ when $X^*=-X$, while differentiating at $t=0$ gives $X^*+X=0$.
So $\dim SO(3)=3$, $\dim SU(2)=3$, $\dim U(n)=n^2$, $\dim SL_2(\R)=3$.

\begin{theorem}[von Neumann, Cartan]\label{thm:closedsubgroup}
Let $G\subset GL_n$ be a matrix Lie group with Lie algebra $\mathfrak g$. There are a neighbourhood $U$ of $0$ in $M_n$ and a neighbourhood $V$ of $I$ such that $\exp$ maps $U\cap\mathfrak g$
homeomorphically onto $V\cap G$. Consequently $G$ is a submanifold of $M_n$ of dimension $\dim\mathfrak g$, with tangent space $T_IG=\mathfrak g$, and the group operations are smooth.
\end{theorem}

\begin{proof}
Choose a complement $W$ of $\mathfrak g$ in $M_n$. \emph{Claim: if $Y_m\in W\setminus\{0\}$, $Y_m\to0$, then $e^{Y_m}\notin G$ for $m$ large.} Otherwise, passing to a subsequence, $e^{Y_m}\in G$ and $Y_m/\|Y_m\|\to Y$, a
unit vector of $W$ (the unit sphere of $W$ is compact). For $t\in\R$ let $k_m$ be the integer part of $t/\|Y_m\|$. Then $e^{k_mY_m}=(e^{Y_m})^{k_m}\in G$, and $k_mY_m=k_m\|Y_m\|\cdot\frac{Y_m}{\|Y_m\|}\to tY$ because $k_m\|Y_m\|\to t$, so
$e^{tY}\in G$ because $G$ is closed. Thus $Y\in\mathfrak g\cap W=\{0\}$, contradicting $\|Y\|=1$.

The map $\Phi(X,Y)=e^Xe^Y$ from $\mathfrak g\times W$ to $GL_n$ has differential $(X,Y)\mapsto X+Y$ at $0$, an isomorphism onto $M_n$, so it is a diffeomorphism from a neighbourhood $U_1\times U_2$ of $0$
onto a neighbourhood $V$ of $I$. If no neighbourhood of $I$ in $G$ is covered by $\exp(U_1\cap\mathfrak g)$, there are $A_m\in G$, $A_m\to I$, with $A_m=e^{X_m}e^{Y_m}$, $X_m\in\mathfrak g$, $Y_m\in W$, $Y_m\ne0$, both tending to
$0$. Then $e^{Y_m}=e^{-X_m}A_m\in G$, contradicting the claim. So, shrinking $V$, $G\cap V=\Phi((U_1\cap\mathfrak g)\times\{0\})=\exp(U_1\cap\mathfrak g)$: in the coordinates given by $\Phi^{-1}$, the set $G\cap V$ is the
linear subspace $\mathfrak g\times\{0\}$, so $G$ is a submanifold near $I$ (Theorem~\ref{thm:fourdesc}). Near any $g\in G$, left multiplication by $g$, a linear isomorphism of $M_n$ preserving $G$, carries this
picture to $g$. Multiplication and inversion are restrictions of smooth maps of $GL_n$ to a submanifold, hence smooth.
\end{proof}

\begin{beauty}
The hypothesis is only that $G$ is a closed subgroup; no smoothness is assumed. Yet the group is automatically a manifold, because the group structure moves every point to every other:
whatever $G$ looks like near the identity, it looks the same everywhere, and a closed set that looks the same everywhere and contains the curves $e^{tX}$ for a whole vector space of $X$ cannot
be fractal. Hilbert's fifth problem asked how far this goes; Gleason, Montgomery and Zippin showed in 1952 that every locally Euclidean topological group is a Lie group.
\end{beauty}

\section{Homomorphisms and the adjoint representation}

\begin{theorem}\label{thm:homlie}
Let $\Phi\colon G\to H$ be a continuous homomorphism of matrix Lie groups. There is a unique linear map $\varphi\colon\mathfrak g\to\mathfrak h$ with $\Phi(e^X)=e^{\varphi(X)}$ for all $X\in\mathfrak g$. It satisfies
$\varphi([X,Y])=[\varphi(X),\varphi(Y)]$ and $\varphi(gXg^{-1})=\Phi(g)\varphi(X)\Phi(g)^{-1}$, and $\varphi=d\Phi_I$.
\end{theorem}

\begin{proof}
$t\mapsto\Phi(e^{tX})$ is a continuous homomorphism $\R\to H$, so it is $e^{t\varphi(X)}$ for a unique $\varphi(X)\in M_n$ by Proposition~\ref{prop:oneparameter}, and $\varphi(X)\in\mathfrak h$ by definition. Clearly $\varphi(sX)=s\varphi(X)$. By
Lie's product formula and continuity,
\[e^{t\varphi(X+Y)}=\Phi\bigl(\lim(e^{tX/m}e^{tY/m})^m\bigr)=\lim(e^{t\varphi(X)/m}e^{t\varphi(Y)/m})^m=e^{t(\varphi(X)+\varphi(Y))},\]
so $\varphi$ is additive; the same argument with the commutator formula gives $e^{t^2\varphi([X,Y])}=e^{t^2[\varphi(X),\varphi(Y)]}$, and differentiating twice at $t=0$ gives the bracket identity. Similarly
$\Phi(ge^{tX}g^{-1})=\Phi(g)e^{t\varphi(X)}\Phi(g)^{-1}$. Finally $\varphi(X)=\frac{\dd}{\dd t}\Phi(e^{tX})|_{t=0}$, which is $d\Phi_I(X)$ once we know that $\Phi$ is smooth: near $I$, $\Phi=\exp\circ\varphi\circ\log$ in the charts of
Theorem~\ref{thm:closedsubgroup}, and elsewhere $\Phi=L_{\Phi(g)}\circ\Phi\circ L_{g^{-1}}$.
\end{proof}

\begin{proposition}\label{prop:connectedgen}
If $G$ is connected, it is generated by $\exp(\mathfrak g)$, and a continuous homomorphism from $G$ is determined by its differential.
\end{proposition}

\begin{proof}
The subgroup $G_0$ generated by $\exp(\mathfrak g)$ contains a neighbourhood $V$ of $I$ in $G$ (Theorem~\ref{thm:closedsubgroup}), so it contains the neighbourhood $gV$ of each of its points $g$: it is open.
Its complement is a union of cosets $hG_0$, also open. As $G$ is connected, $G_0=G$. A homomorphism is determined on $\exp(\mathfrak g)$ by $\varphi$.
\end{proof}

The group $G$ acts on its own Lie algebra by $\mathrm{Ad}_g(X)=gXg^{-1}$ (Proposition~\ref{prop:liealgebra}). This is a continuous homomorphism $\mathrm{Ad}\colon G\to GL(\mathfrak g)$, the \emph{adjoint representation}, and by
Theorem~\ref{thm:homlie} its differential is $\mathrm{ad}\colon\mathfrak g\to\mathfrak{gl}(\mathfrak g)$, $\mathrm{ad}_X(Y)=\frac{\dd}{\dd t}e^{tX}Ye^{-tX}|_{t=0}=[X,Y]$. In particular $\mathrm{Ad}_{e^X}=e^{\mathrm{ad}_X}$.

\section{Rotations and the spin group}

The Lie algebra $\mathfrak{so}(3)$ consists of the antisymmetric matrices
\[L(u)=\begin{pmatrix}0&-u_3&u_2\\u_3&0&-u_1\\-u_2&u_1&0\end{pmatrix},\qquad L(u)v=u\times v,\]
and $[L(u),L(v)]=L(u\times v)$ (by the Jacobi identity of the cross product). So $\mathfrak{so}(3)$ is $\R^3$ with the cross product. Since $L(u)$ has eigenvalues $0,\pm i|u|$ with kernel $\R u$, $e^{L(u)}$ is the
rotation about $u$ by the angle $|u|$ (on $u^\perp$, $L(u)$ is $|u|$ times a quarter turn, and its exponential is a rotation by $|u|$, as in Tome~S1). By Euler's lemma (Tome~2, Lemma~12.3), $\exp\colon\mathfrak{so}(3)\to SO(3)$ is
onto, and $SO(3)$ is path-connected.

The Lie algebra $\mathfrak{su}(2)$ consists of the traceless antihermitian matrices
\[X(u)=\frac12\begin{pmatrix}-iu_3&-u_2-iu_1\\u_2-iu_1&iu_3\end{pmatrix}=-\frac i2(u_1\sigma_1+u_2\sigma_2+u_3\sigma_3),\]
with the Pauli matrices $\sigma_1=\begin{pmatrix}0&1\\1&0\end{pmatrix}$, $\sigma_2=\begin{pmatrix}0&-i\\i&0\end{pmatrix}$, $\sigma_3=\begin{pmatrix}1&0\\0&-1\end{pmatrix}$. From $\sigma_1\sigma_2=i\sigma_3$ and its cyclic permutations,
$[X(u),X(v)]=X(u\times v)$: \emph{$\mathfrak{su}(2)$ and $\mathfrak{so}(3)$ are isomorphic Lie algebras}. The group $SU(2)$ is
\[SU(2)=\left\{\begin{pmatrix}a&-\bar b\\b&\bar a\end{pmatrix}:|a|^2+|b|^2=1\right\},\]
which is the unit sphere $S^3\subset\C^2$: a column $(a,b)$ of length $1$ determines the matrix. (In Tome~2 these were the unit quaternions.)

\begin{theorem}[the spin double cover]\label{thm:spin}
The adjoint representation $\mathrm{Ad}\colon SU(2)\to GL(\mathfrak{su}(2))$, read in the basis $X(e_1),X(e_2),X(e_3)$, is a continuous surjective homomorphism $\pi\colon SU(2)\to SO(3)$ with kernel $\{\pm I\}$, and its differential is
the isomorphism $X(u)\mapsto L(u)$. It is a two-sheeted covering, and the universal covering of $SO(3)$. Consequently $SO(3)$ is homeomorphic to $\R P^3$ and
\[\pi_1(SO(3))\cong\Z/2.\]
The loop $\theta\mapsto e^{\theta L(u)}$, $0\le\theta\le2\pi$, of a full turn about an axis is not nullhomotopic, and the loop of two full turns is.
\end{theorem}

\begin{proof}
\emph{Into $SO(3)$.} The inner product $\langle X,Y\rangle=-2\operatorname{tr}(XY)$ on $\mathfrak{su}(2)$ makes $X(e_1),X(e_2),X(e_3)$ orthonormal (since $\operatorname{tr}(\sigma_j\sigma_k)=2\delta_{jk}$), and it is invariant:
$\operatorname{tr}(gXg^{-1}gYg^{-1})=\operatorname{tr}(XY)$. So $\pi(g)\in O(3)$; $\det\pi(g)$ is $\pm1$, continuous on the connected $S^3$, and $1$ at $I$: $\pi(g)\in SO(3)$.

\emph{Differential.} By Theorem~\ref{thm:homlie} it is $\mathrm{ad}$: $\mathrm{ad}_{X(u)}X(v)=X(u\times v)$, so $\mathrm{ad}_{X(u)}$ is the matrix $L(u)$. It is an isomorphism $\mathfrak{su}(2)\to\mathfrak{so}(3)$, so
$\pi(e^{X(u)})=e^{L(u)}$ for all $u$, and $\exp\colon\mathfrak{so}(3)\to SO(3)$ is onto: \emph{$\pi$ is surjective}.

\emph{Kernel.} $\pi(g)=I$ means $g$ commutes with all of $\mathfrak{su}(2)$, hence with $\sigma_1,\sigma_2,\sigma_3$ and with all $2\times2$ complex matrices (they are spanned by $I$ and the $\sigma_j$). So $g$ is a scalar of
determinant $1$ in $U(2)$: $g=\pm I$.

\emph{Covering.} $\pi$ is continuous, surjective and two-to-one, with $\pi(g)=\pi(-g)$. It induces a continuous bijection $S^3/\{\pm1\}=\R P^3\to SO(3)$, from a compact space onto a Hausdorff one, hence a
homeomorphism, and $\pi$ becomes the quotient map $S^3\to\R P^3$, a two-sheeted covering (Example~\ref{ex:coverings}(d)). Since $S^3$ is simply connected (Corollary~\ref{cor:spheres}), it is the universal covering,
and $\pi_1(SO(3))\cong\{\pm I\}\cong\Z/2$ (Example~\ref{ex:computations}(b)).

\emph{The loops.} Take $u$ of length $1$. The loop $\theta\mapsto e^{\theta L(u)}$ lifts to $\theta\mapsto e^{\theta X(u)}$, which runs from $I$ to $e^{2\pi X(u)}$. Since $X(u)$ has eigenvalues $\pm\frac i2$, $e^{2\pi X(u)}$ has
eigenvalues $e^{\pm i\pi}=-1$: it is $-I$. The lift does not close, so the loop is not nullhomotopic (Proposition~\ref{prop:fibre}); two turns lift to a path from $I$ to $e^{4\pi X(u)}=I$, a closed loop, so
they are nullhomotopic.
\end{proof}

\begin{figure}[ht]
\centering
\fig{ch07-ballmodel}
\caption{$SO(3)$ as a ball of radius $\pi$: the point $\theta u$ is the rotation by $\theta$ about $u$, and antipodal points of the boundary sphere are glued (a rotation by $\pi$ about $u$ equals one about
$-u$). Left: a full turn about an axis is a diameter, a closed loop that cannot be shrunk; its two ends are glued, and moving one end drags the other to the opposite side. Right: two turns are two
diameters; rotating the second one until it meets the first, then pulling the pair inwards, shrinks the loop to a point. This is Dirac's belt trick.}
\label{fig:ballmodel}
\end{figure}

\begin{remark}
The same happens in every dimension $n\ge3$: $\pi_1(SO(n))=\Z/2$, and the universal cover is the double cover $\mathrm{Spin}(n)\to SO(n)$, constructed with Clifford algebras. In the plane
$\pi_1(SO(2))=\pi_1(S^1)=\Z$: a turn by $4\pi$ in the plane cannot be undone, which is why the belt trick needs a third dimension.
\end{remark}

\section{$SL_2(\R)$ and the maximal torus of $U(n)$}

The Lie algebra $\mathfrak{sl}_2(\R)$ has the basis
\[H=\begin{pmatrix}1&0\\0&-1\end{pmatrix},\quad E=\begin{pmatrix}0&1\\0&0\end{pmatrix},\quad F=\begin{pmatrix}0&0\\1&0\end{pmatrix},\qquad[H,E]=2E,\ [H,F]=-2F,\ [E,F]=H.\]
It is not isomorphic to $\mathfrak{so}(3)$: in $\mathfrak{so}(3)$ no $\mathrm{ad}_X$ has a nonzero real eigenvalue ($\mathrm{ad}_{L(u)}=L(u)$ has eigenvalues $0,\pm i|u|$), while $\mathrm{ad}_H$ has eigenvalue $2$ on $E$. The three kinds of
one-parameter subgroups, $e^{tH}$ (hyperbolic), $e^{tE}$ (parabolic, unipotent) and rotations $e^{t(E-F)}$ (elliptic), will be the three kinds of isometries of the hyperbolic plane in Chapter~9. Unlike the compact
groups, $SL_2(\R)$ is not the image of its exponential (Exercise~\ref{ex:slexp}).

\begin{theorem}[maximal torus and Weyl group of $U(n)$]\label{thm:torus}
Let $T\subset U(n)$ be the subgroup of diagonal unitary matrices, a torus $(S^1)^n$.
\begin{enumerate}
\item[(a)] Every $g\in U(n)$ is conjugate in $U(n)$ to an element of $T$; consequently $\exp\colon\mathfrak u(n)\to U(n)$ is onto, and $U(n)$ is connected.
\item[(b)] If $t\in T$ has distinct diagonal entries, the matrices commuting with $t$ in $U(n)$ are exactly $T$; so $T$ is a maximal abelian subgroup.
\item[(c)] The normalizer $N(T)=\{g:gTg^{-1}=T\}$ consists of the matrices $P_wD$, with $D\in T$ and $P_w$ the permutation matrix of $w\in S_n$, and the \emph{Weyl group} $N(T)/T$ is isomorphic to $S_n$, acting on $T$ by
permuting the diagonal entries.
\item[(d)] Two elements of $T$ are conjugate in $U(n)$ if and only if they differ by a permutation of the diagonal entries.
\end{enumerate}
\end{theorem}

\begin{proof}
(a) A unitary matrix is normal, so it is $UDU^*$ with $U$ unitary and $D$ diagonal (Tome~S1, Theorem~5.2); the eigenvalues, being of modulus $1$, are $e^{i\theta_j}$, so $D=e^{i\Theta}$ with $\Theta$ real diagonal and
$g=e^{U(i\Theta)U^*}$ with $U(i\Theta)U^*\in\mathfrak u(n)$. The path $s\mapsto e^{sU(i\Theta)U^*}$ joins $I$ to $g$. (b) If $gt=tg$, $g$ preserves each eigenspace of $t$, which are the coordinate lines, so $g$ is
diagonal. (c) Fix $t\in T$ with distinct entries. If $g\in N(T)$, then $gtg^{-1}\in T$ has the same eigenvalues as $t$, so $gtg^{-1}=P_wtP_w^{-1}$ for a permutation $w$, and $P_w^{-1}g$ commutes with $t$: by (b), $g=P_wD$.
Conversely $P_wDtD^{-1}P_w^{-1}=P_wtP_w^{-1}$, which permutes the entries of $t$. The map $P_wD\mapsto w$ is a homomorphism onto $S_n$ with kernel $T$. (d) Conjugate matrices have the same eigenvalues with
multiplicity, and permutations are realized by the $P_w$.
\end{proof}

So conjugacy classes of $U(n)$ correspond to unordered $n$-tuples of points on the unit circle: the eigenvalues. Chapter~8 integrates over $U(n)$ by integrating over the torus with the right weight,
and the weight comes from the Weyl group and the adjoint representation.

\section{Exercises}

\begin{exercise}[$\star\star$]\label{ex:slexp}
\emph{The exponential of $SL_2(\R)$ is not onto.} (a) Show that if $X\in\mathfrak{sl}_2(\R)$, then $X^2=-\det(X)\,I$, and compute $e^X$: it is $\cosh\sqrt{-\det X}\,I+\frac{\sinh\sqrt{-\det X}}{\sqrt{-\det X}}X$ when $\det X<0$, with the obvious
variants when $\det X\ge0$. (b) Deduce that $\operatorname{tr}e^X\ge-2$, with equality only for $e^X=-I$. (c) Show that $\begin{pmatrix}-2&0\\0&-\frac12\end{pmatrix}$ and $\begin{pmatrix}-1&1\\0&-1\end{pmatrix}$ are not exponentials, but are products of two exponentials.
\end{exercise}

\begin{exercise}[$\star$]\label{ex:heisenberg}
\emph{The Heisenberg group.} Let $G$ be the group of matrices $\begin{pmatrix}1&a&c\\0&1&b\\0&0&1\end{pmatrix}$. (a) Find its Lie algebra $\mathfrak h$ and its brackets; show that $[\mathfrak h,\mathfrak h]$ is central. (b) Show that $\exp\colon\mathfrak h\to G$ is a
bijection and that $e^Xe^Y=e^{X+Y+\frac12[X,Y]}$. (c) In quantum mechanics position and momentum satisfy $[Q,P]=i\hbar$; relate this to (a).
\end{exercise}

\begin{exercise}[$\star$]\label{ex:oneparamso3}
\emph{One-parameter subgroups of $SO(3)$.} Show that every continuous homomorphism $\R\to SO(3)$ is $t\mapsto$ rotation about a fixed axis $u$ by the angle $\omega t$, for some $u$ and $\omega$. Which of them are periodic?
Show that two rotations commute if and only if they have the same axis or are two half-turns about perpendicular axes.
\end{exercise}

\begin{exercise}[$\star\star$]\label{ex:sondconnected}
\emph{$SO(n)$.} (a) Show that every $A\in SO(n)$ is conjugate in $SO(n)$ to a block diagonal matrix of plane rotations and $\pm1$'s with an even number of $-1$'s, and deduce that $\exp\colon\mathfrak{so}(n)\to SO(n)$ is onto and $SO(n)$ is
connected. (b) Show that $O(n)$ has exactly two connected components.
\end{exercise}

\begin{exercise}[$\star\star$]\label{ex:sl2pi1}
\emph{The fundamental group of $SL_2(\R)$.} Using the polar decomposition $A=WP$ (Tome~S1, Theorem~6.3), show that $SL_2(\R)$ is homeomorphic to $SO(2)\times\{P\text{ symmetric positive definite},\ \det P=1\}\cong S^1\times\R^2$, and deduce
$\pi_1(SL_2(\R))\cong\Z$. (Unlike $SO(3)$, its universal cover is not a matrix group; this remark is for culture.)
\end{exercise}

\begin{exercise}[$\star$]\label{ex:centersu2}
\emph{Centres.} Show that the centre of $SU(2)$ is $\{\pm I\}$ and the centre of $SO(3)$ is trivial; deduce again that $SU(2)$ and $SO(3)$ are not isomorphic as groups, although their Lie algebras are.
\end{exercise}

\section{Hints}

\begin{hint}{ex:slexp}
(a) Cayley--Hamilton for a traceless $2\times2$ matrix. (b) $\operatorname{tr}e^X$ is $2\cosh\lambda$, $2$ or $2\cos\mu$. (c) For the products, write a negative diagonal matrix as $-I$ times a positive one.
\end{hint}

\begin{hint}{ex:heisenberg}
$\mathfrak h$ is the strictly upper triangular matrices; products of three of them vanish. For (b), use the series of $e^X$, which stops at $X^2$.
\end{hint}

\begin{hint}{ex:oneparamso3}
Proposition~\ref{prop:oneparameter} and the description of $e^{L(u)}$. For commuting rotations, a rotation commuting with $R$ maps the axis of $R$ to itself.
\end{hint}

\begin{hint}{ex:sondconnected}
Use the real normal form of an orthogonal matrix: complex eigenvalues come in conjugate pairs, each giving a plane rotation. Pair two $-1$'s into a rotation by $\pi$.
\end{hint}

\begin{hint}{ex:sl2pi1}
The polar factors depend continuously on $A$, since $P=\sqrt{A^{\mathsf T}A}$ and $W=AP^{-1}$. A positive definite symmetric matrix of determinant $1$ is $e^S$ with $S$ symmetric of trace $0$.
\end{hint}

\begin{hint}{ex:centersu2}
A central element of a connected matrix group commutes with $e^{tX}$ for all $X$, hence with all of $\mathfrak g$.
\end{hint}

\section{Solutions}

\begin{solution}{ex:slexp}
\idea Every element of $\mathfrak{sl}_2$ squares to a scalar, so its exponential is a combination of $I$ and $X$, and its trace is a $\cosh$, $2$ or a $\cos$: never less than $-2$.
\writeup (a) $X$ has trace $0$, so by Cayley--Hamilton $X^2=-\det(X)I$. With $\lambda=\sqrt{-\det X}$ when $\det X<0$, the even powers are $\lambda^{2k}I$ and the odd ones $\lambda^{2k}X$, so $e^X=\cosh\lambda\,I+\frac{\sinh\lambda}\lambda X$. When
$\det X=\mu^2>0$, $e^X=\cos\mu\,I+\frac{\sin\mu}\mu X$, and when $\det X=0$, $e^X=I+X$. (b) The traces are $2\cosh\lambda\ge2$, $2\cos\mu\ge-2$ and $2$. If $\operatorname{tr}e^X=-2$, then $\cos\mu=-1$, $\sin\mu=0$ and $e^X=-I$.
(c) The first matrix has trace $-\frac52<-2$; the second has trace $-2$ and is not $-I$. But $\begin{pmatrix}-2&0\\0&-\frac12\end{pmatrix}=e^{\pi(E-F)}e^{(\log2)H}$, since $e^{\pi(E-F)}=-I$, and
$\begin{pmatrix}-1&1\\0&-1\end{pmatrix}=e^{\pi(E-F)}e^{-E}$. In accordance with Proposition~\ref{prop:connectedgen}, $SL_2(\R)$ (which is connected, Exercise~\ref{ex:sl2pi1}) is generated by exponentials, but needs products.
\end{solution}

\begin{solution}{ex:heisenberg}
\idea Nilpotent matrices have finite exponential series, so everything can be computed by hand.
\writeup (a) $e^{tX}$ is unipotent upper triangular for all $t$ exactly when $X$ is strictly upper triangular: $\mathfrak h$ is spanned by $P=E_{12}$, $Q=E_{23}$, $Z=E_{13}$, with $[P,Q]=Z$ and $Z$ commuting
with everything. (b) $X^3=0$, so $e^X=I+X+\frac12X^2$, and $X=(x_{12},x_{23},x_{13})$ gives $e^X$ with entries $a=x_{12}$, $b=x_{23}$, $c=x_{13}+\frac12x_{12}x_{23}$, which can be solved uniquely for $X$: a bijection. Since
$[X,Y]$ is central and $[X,[X,Y]]=0$, $e^Xe^Y=I+X+Y+\frac{X^2}2+XY+\frac{Y^2}2$ (all products of three vanish), and $e^{X+Y+\frac12[X,Y]}=I+X+Y+\frac12[X,Y]+\frac12(X+Y)^2$, which is the same since
$\frac12(XY-YX)+\frac12(XY+YX)=XY$. (c) $[Q,P]$ is a multiple of the identity, which commutes with everything: the canonical commutation relations define a representation of $\mathfrak h$ (with $Z$ acting
by $i\hbar$), and the Heisenberg group is named after them.
\end{solution}

\begin{solution}{ex:oneparamso3}
\idea A one-parameter subgroup is $e^{tX}$, and $e^{tL(u)}$ turns about $u$ at constant speed.
\writeup By Proposition~\ref{prop:oneparameter} the homomorphism is $e^{tL(v)}$ for some $v$; with $u=\frac v{|v|}$ and $\omega=|v|$ (any $u$ if $v=0$), $e^{tL(v)}$ is the rotation by $\omega t$ about $u$. All are periodic, of period $\frac{2\pi}\omega$
(or constant). If $RS=SR$ and $R$ is not the identity, $S$ maps the axis of $R$ (its fixed line, if $R$ is not a half-turn; the $+1$ eigenspace in general) to itself, so $Su=\pm u$. If $Su=u$, $S$ has the same axis. If
$Su=-u$, $S$ is a half-turn about an axis $w\perp u$, and $SRS^{-1}$ is the rotation by the same angle about $Su=-u$, that is the rotation by $-\theta$ about $u$; it equals $R$ only if $\theta=\pi$. Conversely, rotations
about the same axis commute, and two half-turns about perpendicular axes commute (their product is the half-turn about the third axis, in either order).
\end{solution}

\begin{solution}{ex:sondconnected}
\idea An orthogonal matrix is a product of plane rotations in perpendicular planes and reflections; in $SO(n)$ the reflections pair up into rotations by $\pi$.
\writeup (a) As a complex matrix, $A$ is unitary, so diagonalizable with eigenvalues of modulus $1$; nonreal ones come in conjugate pairs $e^{\pm i\theta}$, with eigenvectors $v,\bar v$, and $\operatorname{Re}v$, $\operatorname{Im}v$ span a real
plane on which $A$ is a rotation by $\theta$. These planes and the real eigenvectors (eigenvalues $\pm1$) are mutually orthogonal (eigenspaces of a normal matrix are), so in a suitable orthonormal basis $A$ is block
diagonal. The determinant is $(-1)^{\#(-1)}$, so the number of $-1$'s is even, and pairs of them are rotations by $\pi$. Changing the sign of one basis vector if needed puts the basis change in $SO(n)$. Each block
is $e^{\theta J}$ with $J$ the quarter turn, so $A=e^X$ with $X$ antisymmetric, and $s\mapsto e^{sX}$ joins $I$ to $A$. (b) $\det$ is continuous with values $\pm1$, so $O(n)$ has at least two components; multiplication by a
reflection maps $SO(n)$ homeomorphically onto $\det=-1$, which is therefore connected.
\end{solution}

\begin{solution}{ex:sl2pi1}
\idea Every matrix of determinant $1$ is a rotation times a stretch, uniquely, and the stretches form a contractible set.
\writeup By the polar decomposition, $A=WP$ with $P=\sqrt{A^{\mathsf T}A}$ positive definite and $W=AP^{-1}$ orthogonal; $\det P=|\det A|=1$, so $\det W=1$. The maps $A\mapsto(W,P)$ and $(W,P)\mapsto WP$ are continuous and
inverse to each other (the square root of a positive definite matrix is continuous, for instance because $\sqrt{M}=\frac1{2\pi i}\oint\sqrt z(z-M)^{-1}\dd z$ on a fixed contour around the spectrum, or by
diagonalization and the uniqueness of $P$). The positive definite symmetric matrices of determinant $1$ are the $e^S$ with $S$ symmetric of trace $0$, that is $S=\begin{pmatrix}x&y\\y&-x\end{pmatrix}$, and $S\mapsto e^S$ is a
homeomorphism from $\R^2$ (its inverse is the logarithm of a positive matrix, computed in an eigenbasis). So $SL_2(\R)\cong S^1\times\R^2$, which deformation retracts onto $S^1$: $\pi_1\cong\Z$, generated by the loop of
rotations $e^{\theta(E-F)}$, $0\le\theta\le2\pi$.
\end{solution}

\begin{solution}{ex:centersu2}
\idea Being central is a linear condition on the Lie algebra, and the Lie algebras $\mathfrak{su}(2)$ and $\mathfrak{so}(3)$ have trivial centre.
\writeup A central $g$ commutes with $e^{tX}$, and differentiating at $t=0$, with every $X\in\mathfrak g$. In $SU(2)$, commuting with the Pauli matrices forces $g$ to be a scalar, so $g=\pm I$. In $SO(3)$, if $g$ commutes with the rotation $R$
about $u$ by a small angle, then $gRg^{-1}$, the rotation about $gu$ by the same angle, equals $R$, so $gu=\pm u$. Thus every vector is an eigenvector of $g$, which is therefore $\pm I$, and $\det g=1$ gives $g=I$. A group isomorphism
maps the centre onto the centre, and $\{\pm I\}$ has two elements while the centre of $SO(3)$ has one.
\end{solution}

\section*{Sources and further reading}
\addcontentsline{toc}{section}{Sources and further reading}

G.~E.~Uhlenbeck and S.~Goudsmit, ``Ersetzung der Hypothese vom unmechanischen Zwang durch eine Forderung bez\"uglich des inneren Verhaltens jedes einzelnen Elektrons'', \emph{Naturwissenschaften}
13 (1925), 953--954; G.~E.~Uhlenbeck, ``Personal reminiscences'', \emph{Physics Today} 29, no.~6 (1976), 43--48, which tells the story of Lorentz's calculation and Ehrenfest's answer; W.~Pauli, ``Zur Quantenmechanik
des magnetischen Elektrons'', \emph{Z.\ Phys.} 43 (1927), 601--623. The neutron experiments: H.~Rauch et al., ``Verification of coherent spinor rotation of fermions'', \emph{Phys.\ Lett.\ A} 54 (1975), 425--427, and
S.~A.~Werner et al., ``Observation of the phase shift of a neutron due to precession in a magnetic field'', \emph{Phys.\ Rev.\ Lett.} 35 (1975), 1053--1055.

J.~von Neumann, ``\"Uber die analytischen Eigenschaften von Gruppen linearer Transformationen und ihrer Darstellungen'', \emph{Math.\ Z.} 30 (1929), 3--42; \'E.~Cartan, \emph{La th\'eorie des groupes finis et
continus et l'Analysis situs} (M\'emorial des sciences math\'ematiques 42, Gauthier-Villars, 1930). A.~Gleason, ``Groups without small subgroups'', \emph{Ann.\ of Math.} 56 (1952), 193--212, and D.~Montgomery and
L.~Zippin, ``Small subgroups of finite-dimensional groups'', same volume, 213--241.

The treatment follows B.~Hall, \emph{Lie Groups, Lie Algebras, and Representations}, 2nd ed. (Springer, 2015), Chapters~1--3 and~11, and W.~Rossmann, \emph{Lie Groups: An Introduction through Linear Groups}
(Oxford University Press, 2002). For the belt trick and its physics, R.~Penrose and W.~Rindler, \emph{Spinors and Space-Time}, vol.~1 (Cambridge University Press, 1984), Chapter~1.
